% =====================================================================
%  PRÉAMBULE COMMUN — Carnet d'entraînement, Fiche 2 (Nomenclature)
%  (inséré physiquement dans chaque document pour qu'il soit autonome)
% =====================================================================
\documentclass[11pt,a4paper]{article}

% ---------- Encodage & langue ----------
\usepackage[utf8]{inputenc}
\usepackage[T1]{fontenc}
\usepackage{lmodern}
\usepackage[french]{babel}

% ---------- Mathématiques & chimie ----------
\usepackage{amsmath,amssymb,mathtools}
\usepackage{icomma}
\usepackage[version=4]{mhchem}
\usepackage{siunitx}
\sisetup{output-decimal-marker={,},per-mode=symbol}

% ---------- Graphismes & mise en page ----------
\usepackage{xcolor}
\usepackage{colortbl}
\usepackage{tikz}
\usepackage[most]{tcolorbox}
\usepackage{enumitem}
\usepackage{array}
\usepackage{booktabs}
\usepackage{multicol}
\usepackage{fancyhdr}
\usepackage{titlesec}
\usepackage[hidelinks]{hyperref}
\usetikzlibrary{arrows.meta,positioning,calc,shapes.geometric,shapes.misc,
  fit,backgrounds,decorations.pathreplacing}
\usepackage[a4paper,margin=2.2cm,headheight=28pt,headsep=10pt]{geometry}

% =====================================================================
%  PALETTE (charte « Chimie » — mauve/prune du cours)
% =====================================================================
\definecolor{primary}{RGB}{146,95,161}
\definecolor{primarydark}{RGB}{92,55,108}
\definecolor{defcol}{RGB}{0,114,178}
\definecolor{thmcol}{RGB}{0,140,120}
\definecolor{formulacol}{RGB}{198,93,9}
\definecolor{remarkcol}{RGB}{120,94,200}
\definecolor{attncol}{RGB}{200,40,55}
\definecolor{excol}{RGB}{0,135,90}
\definecolor{methcol}{RGB}{90,90,100}
\definecolor{metalcol}{RGB}{198,150,40}
\definecolor{nonmetcol}{RGB}{0,140,150}
\definecolor{oxycol}{RGB}{210,55,45}
\definecolor{hydcol}{RGB}{60,120,210}
\definecolor{catcol}{RGB}{200,70,120}
\definecolor{ancol}{RGB}{40,110,190}
\definecolor{saltcol}{RGB}{120,90,160}

% =====================================================================
%  STYLE COMMUN DES BOÎTES
% =====================================================================
\tcbset{
  boxbase/.style={enhanced,breakable,boxrule=0.9pt,arc=2.5pt,
     left=3mm,right=3mm,top=2.3mm,bottom=2.3mm,
     fonttitle=\bfseries,coltitle=white,
     toptitle=0.6mm,bottomtitle=0.6mm},
}

\newtcolorbox{definition}[1][]{%
  boxbase,colback=defcol!5,colframe=defcol,
  title={\textsc{Définition}\ifx\relax#1\relax\relax\else\textmd{ --- \itshape #1}\fi}}
\newtcolorbox{regle}[1][]{%
  boxbase,colback=thmcol!6,colframe=thmcol,
  title={\textsc{Règle}\ifx\relax#1\relax\relax\else\textmd{ --- \itshape #1}\fi}}
\newtcolorbox{formule}[1][]{%
  boxbase,colback=formulacol!6,colframe=formulacol,
  title={\textsc{Méthode-clé}\ifx\relax#1\relax\relax\else\textmd{ --- \itshape #1}\fi}}
\newtcolorbox{remarque}[1][]{%
  boxbase,colback=remarkcol!6,colframe=remarkcol,
  title={\textsc{Remarque}\ifx\relax#1\relax\relax\else\textmd{ : \itshape #1}\fi}}
\newtcolorbox{attention}[1][]{%
  boxbase,colback=attncol!6,colframe=attncol,
  title={\textsc{Attention}\ifx\relax#1\relax\relax\else\textmd{ : \itshape #1}\fi}}
\newtcolorbox{exemple}[1][]{%
  boxbase,colback=excol!5,colframe=excol,
  title={\textsc{Exemple}\ifx\relax#1\relax\relax\else\textmd{ : \itshape #1}\fi}}
\newtcolorbox{methode}[1][]{%
  boxbase,colback=methcol!7,colframe=methcol,sharp corners,
  title={\textsc{Méthode}\ifx\relax#1\relax\relax\else\textmd{ : \itshape #1}\fi}}
\newtcolorbox{astuce}[1][]{%
  boxbase,colback=primary!6,colframe=primary,
  title={\textsc{Astuce concours}\ifx\relax#1\relax\relax\else\textmd{ : \itshape #1}\fi}}

% ---- Exercice (numéroté) ----
\newtcolorbox[auto counter]{exercice}[1][]{%
  boxbase,colback=primary!4,colframe=primarydark,
  title={\textsc{Exercice}~\thetcbcounter\ifx\relax#1\relax\relax\else\textmd{ --- \itshape #1}\fi}}

% ---- Corrigé (numéroté) ----
\newtcolorbox[auto counter]{corrige}[1][]{%
  boxbase,colback=excol!4,colframe=excol,
  title={\textsc{Corrigé}~\thetcbcounter\ifx\relax#1\relax\relax\else\textmd{ --- \itshape #1}\fi}}

% =====================================================================
%  EN-TÊTES ET PIEDS DE PAGE
% =====================================================================
\pagestyle{fancy}
\fancyhf{}
\fancyhead[L]{\small\textcolor{primarydark}{\textbf{Fiche 2} — Nomenclature}}
\fancyhead[R]{\small\textcolor{primarydark}{\textsc{Carnet d'entraînement}}}
\fancyfoot[C]{\small\thepage}
\renewcommand{\headrulewidth}{0.8pt}
\renewcommand{\headrule}{\hbox to\headwidth{\color{primary}\leaders\hrule height \headrulewidth\hfill}}

\titleformat{\section}{\Large\bfseries\color{primarydark}}{\thesection}{0.6em}{}
\titleformat{\subsection}{\large\bfseries\color{primary}}{\thesubsection}{0.6em}{}
\titleformat{\subsubsection}{\normalsize\bfseries\color{primary!80!black}}{}{0em}{}

% ---- Bandeau de titre réutilisable : \bandeau{Thème}{Sous-titre} ----
\newcommand{\bandeau}[2]{%
  \begin{center}
  \begin{tikzpicture}
  \node[rectangle,rounded corners=7pt,fill=primary,text=white,
        minimum width=16.0cm,minimum height=1.7cm,align=center,inner sep=8pt]
    {{\Large\bfseries #1}\\[3pt]{\normalsize #2}};
  \end{tikzpicture}
  \end{center}
  \vspace{0.3cm}}
\begin{document}
\bandeau{Thème 1 — Difficile}{Ions, valences et nombre d'oxydation --- Corrigé}

\begin{corrige}[le permanganate de potassium, décortiqué]
\begin{enumerate}[label=\alph*),leftmargin=1.8em,topsep=2pt,itemsep=1.5pt]
  \item \ce{K} est un alcalin $\Rightarrow +\mathrm{I}$ (convention 5) ; \ce{O} $\Rightarrow -\mathrm{II}$ (convention 3).
  \item Neutralité : $(+\mathrm{I}) + x + 4(-\mathrm{II}) = 0 \Rightarrow x = 8-1 = +\mathrm{VII}$.
  \item Le manganèse ($Z=25$) a cédé 7 électrons : il en reste $25 - 7 = 18$.
  \item \textbf{Faux.} Le calcul donne $+\mathrm{VII}$, pas $+\mathrm{VIII}$. L'erreur classique consiste à
        oublier le $+\mathrm{I}$ du potassium (on écrirait alors $x = 8$). Le degré maximal réel du manganèse
        est bien $+\mathrm{VII}$ (celui du permanganate).
  \item \ce{K}\,($+\mathrm{I}$)\ \ce{Mn}\,($+\mathrm{VII}$)\ \ce{O}\,($-\mathrm{II}$).
\end{enumerate}
\end{corrige}

\begin{corrige}[l'oxyde de cuivre \ce{Cu2O}]
\begin{enumerate}[label=\alph*),leftmargin=1.8em,topsep=2pt,itemsep=1.5pt]
  \item \ce{Cu2O} neutre, \ce{O}$=-\mathrm{II}$ : $2x + (-2) = 0 \Rightarrow x = +\mathrm{I}$. Cuivre au degré $+\mathrm{I}$.
  \item Cuivre : $29 - 1 = 28$ électrons. Oxygène (\ce{O^2-}, a gagné 2) : $8 + 2 = 10$ électrons.
  \item Le cuivre est un \textbf{métal} : à l'état naturel son n.o.\ ne peut pas être négatif (un métal
        \emph{cède} des électrons). « Cuivre $-\mathrm{I}$ » est donc physiquement impossible.
  \item \ce{Cu2O} = \textbf{oxyde de cuivre(I)} (ou oxyde cuivreux). \ce{CuO} est l'\textbf{oxyde de cuivre(II)}
        (cuivrique) : même couple d'éléments, mais le degré du cuivre change (I vs II), donc la formule change.
\end{enumerate}
\end{corrige}

\begin{corrige}[quand le n.o.\ est une moyenne : le thiosulfate]
\begin{enumerate}[label=\alph*),leftmargin=1.8em,topsep=2pt,itemsep=1.5pt]
  \item \ce{SO4^2-} : $x + 4(-2) = -2 \Rightarrow x = +\mathrm{VI}$.
  \item \ce{S2O3^2-} : $2x + 3(-2) = -2 \Rightarrow 2x = 4 \Rightarrow x = +\mathrm{II}$ (moyenne). C'est bien
        inférieur au $+\mathrm{VI}$ du sulfate : le thiosulfate contient un soufre « supplémentaire » à la
        place d'un oxygène, ce qui abaisse le degré moyen.
  \item La méthode « somme des n.o.\ $=$ charge » répartit également la charge sur les 2 soufres : elle donne
        une \emph{moyenne}. En réalité, les deux atomes de soufre du thiosulfate ne sont pas dans le même
        environnement (l'un est central, l'autre remplace un oxygène), donc leurs degrés réels diffèrent ;
        seule la moyenne vaut $+\mathrm{II}$.
\end{enumerate}
\end{corrige}

\begin{corrige}[la réaction aluminothermique : transferts d'électrons]
\begin{enumerate}[label=\alph*),leftmargin=1.8em,topsep=2pt,itemsep=1.5pt]
  \item \ce{Fe2O3} : $2x + 3(-2)=0 \Rightarrow x = +\mathrm{III}$. Dans \ce{Fe} (métal) : $0$.
        Le fer passe de $+\mathrm{III}$ à $0$ : il est \emph{réduit} (il capte 3 électrons par atome).
  \item \ce{Al} (métal) : $0$. \ce{Al2O3} : $2x + 3(-2)=0 \Rightarrow x=+\mathrm{III}$.
        L'aluminium passe de $0$ à $+\mathrm{III}$ : il est \emph{oxydé} (c'est le réducteur).
  \item Chaque atome d'aluminium cède \textbf{3 électrons} ($0 \to +\mathrm{III}$).
  \item Bilan : 2 \ce{Al} cèdent $2\times3 = 6$ électrons ; 2 \ce{Fe} captent $2\times3 = 6$ électrons.
        Les 6 électrons cédés sont exactement captés : le bilan est équilibré.
\end{enumerate}
\end{corrige}

\begin{corrige}[les exceptions : peroxyde et hydrure]
\begin{enumerate}[label=\alph*),leftmargin=1.8em,topsep=2pt,itemsep=1.5pt]
  \item \ce{H2O2} : $2(+1) + 2x = 0 \Rightarrow 2x = -2 \Rightarrow x = -\mathrm{I}$. Ce n'est \emph{pas} la
        valeur habituelle $-\mathrm{II}$ : dans un \textbf{peroxyde}, l'oxygène vaut $-\mathrm{I}$ (exception
        de la convention 3, à cause de la liaison \ce{O-O}).
  \item \ce{NaH} : $(+\mathrm{I}) + x = 0 \Rightarrow x = -\mathrm{I}$. L'hydrogène est ici à $-\mathrm{I}$
        (ion hydrure \ce{H-}), ce qui est inhabituel : d'ordinaire \ce{H} vaut $+\mathrm{I}$.
  \item Les conventions sont mises en défaut : (i) pour l'\textbf{oxygène}, dans les \emph{peroxydes}
        ($-\mathrm{I}$) ; (ii) pour l'\textbf{hydrogène}, dans les \emph{hydrures métalliques} ($-\mathrm{I}$).
\end{enumerate}
\end{corrige}
\end{document}
